% !TEX TS-program = xelatex
% !TEX encoding = UTF-8 Unicode
% !Mode:: "TeX:UTF-8"

\documentclass{my-resume}
\usepackage{my-resume-zh} % Simplified Chinese support (HarmonyOS Sans SC, self-contained)
\usepackage{linespacing_fix} % disable extra space before next section
\usepackage{cite}

\begin{document}
\pagenumbering{gobble} % suppress displaying page number

\photoheader{%
  \name{吴筱琦}
  \contactInfo{+86 13776020628}{yukiwxq@gmail.com}{现居城市：伦敦/苏州}{意向城市：上海/苏州/杭州 \textbar\ 数据分析与生态环境科学}
}{photo.jpg}{2.6cm}

\section{个人总结}
帝国理工学院计算生命科学研究型硕士在读生，具备扎实的生命科学背景，系统掌握 R、Python、统计建模、生物信息学及基于大语言模型的数据提取流程。独立设计并实现 VecTraits Curation 项目——一套从科学文献中提取结构化性状数据的端到端 LLM 流程；曾在上海宣泰医药担任分析工程师，主导 HPLC/UV 稳定性检测与质量控制，积累医药研发一线经验。在毕业设计及 11 项独立研究项目中，展现出严谨的定量分析能力与出色的科研写作、跨团队协作能力。希望将统计编程与生物学洞见相结合，为医药研发或临床数据岗位提供数据驱动的分析支持。

\section{教育背景}
\entry{帝国理工学院 (QS 2)}{2025.09 - 2026.09}{计算生命科学研究型硕士}{伦敦}
\begin{itemize}[parsep=0.2ex]
  \item \textbf{课程内容}：R语言统计学、生态与进化数据科学、生物计算训练营等生态与进化定量方法核心课程，涵盖统计编程、数据科学、GIS、基因组学与生物信息学
  \item \textbf{独立研究}：开展为期8个月的独立研究项目，聚焦编程、建模与统计分析在生态与进化问题中的应用
\end{itemize}
\entry{帝国理工学院 (QS 2)}{2022.09 - 2025.07}{生物科学 本科}{伦敦}
\begin{itemize}[parsep=0.2ex]
  \item \textbf{相关课程}：进化与多样性，生物化学与微生物学，生态学与进化论，细胞生物学与遗传学，应用分子生物学，遗传学，生态学基础，行为生态学，生态学实地技能，生物信息学，编程与统计，非洲生物学野外课程，生物多样性与保护生物学，全球变化生物学，法语
\end{itemize}

\section{实习与社会实践}
\entry{上海宣泰医药科技有限公司}{2024.07 - 2024.09}{分析工程师 \textperiodcentered\ 分析部}{上海}
\begin{itemize}[parsep=0.2ex]
  \item \textbf{稳定性检测}：使用高效液相色谱-紫外检测（HPLC/UV）技术对客户提供的药品样本进行稳定性检测，专注于样品中杂质的检测与分析
  \item \textbf{数据分析}：独立完成样本前处理、HPLC系统操作及数据分析，确保检测结果的准确性和可靠性
  \item \textbf{质量控制}：通过严格的质控流程鉴别潜在杂质，为药品的稳定性评估提供关键数据支持
\end{itemize}
\entry{上海剑藤教育科技有限公司}{2023.07 - 2023.08}{教学助教 \textperiodcentered\ 生化部}{上海}
\begin{itemize}[parsep=0.2ex]
  \item \textbf{教学辅导}：为冲刺牛津、剑桥等名校的学生提供课外拓展教学辅导，受到国际高中孩子们的欢迎
  \item \textbf{教学执行}：利用本科专业知识，承担课程的部分全英文讲授、批改作业、作业讲解、模拟面试、组织课堂讨论等教学工作，保证教学质量
  \item \textbf{教学成果}：整理和分析输出需求文档，帮助老师根据学生的特点，差生达到中等成绩、优等生进一步巩固基础，最终帮助多名学生拿到了剑桥本科录取通知
\end{itemize}

\section{科研与项目经历}

\entry{多模态大语言模型提取媒介性状数据评估：高精确率与结构依赖型记录召回}{2025.12 - 2026.08}{}{伦敦}
\role{Python、Gemma 4 31B IT（多模态LLM）、PyMuPDF、Docling、提示工程、JSON Schema 校验、二分图记录匹配评估、pytest}{}
\begin{itemize}[parsep=0.2ex]
  \item \textbf{基准构建}：从 VecTraits 平台超4万行、334篇论文的原始导出数据中，设计加权多样性贪心算法选取20篇论文（覆盖不同性状类型、报告格式与数据来源），构建含2,335条人工基准记录的开发集与431条记录的留出测试集；参照 MIReVTD 标准将91字段的VecTraits原始schema精简为22字段提取schema
  \item \textbf{多模态提取流程}：基于 PyMuPDF 提取全文文本、Docling 解析表格结构、150 DPI 渲染图表证据构建多模态证据包，输入 Gemma 4 31B IT（256K 上下文多模态大语言模型），通过提示工程与 JSON Schema 约束实现从 PDF 文本/表格/图像中可溯源提取原始性状记录；经11轮提示词迭代优化，最终版本在开发集上取得 0.810 的条件宏观字段 F1
  \item \textbf{严谨评估与结果}：设计基于二分图记录匹配的评估框架（行集合/条件字段/端到端字段三类指标）及错误溯源分析；在留出测试集上评估锁定流程，取得记录精确率 1.000、召回率 0.360、条件宏观字段 F1 0.729，验证流程"高精确率、结构依赖型召回"特性，论证其适合作为人机协同数据整理的候选记录生成工具
\end{itemize}

\entry{气候变化对佛罗里达州埃及伊蚊种群动态的影响建模}{2025.03 - 2025.06}{}{伦敦}
\role{R、deSolve、ggplot2、ODEs、统计验证（RMSE、Pearson、Spearman）}{}
\begin{itemize}[parsep=0.2ex]
  \item \textbf{模型适配}：将已发表的热带 ODE 房室模型适配至亚热带地理环境，模拟佛罗里达 Collier 县埃及伊蚊种群在温度与降雨影响下的动态变化
  \item \textbf{模型验证}：基于 53 周本地蚊虫监测数据验证模型（RMSE=1.65，Pearson r=0.29，Spearman $\rho$=0.30），优化时间滞后（$\tau=-2$ 周）与降雨尺度因子，量化模型性能
  \item \textbf{研究成果}：完成毕业设计，提出模型改进方向（引入人为水源、空间异质性与不确定性量化），为媒介控制与疾病防控研究提供参考
\end{itemize}

\entry{北大西洋渔业数据变化驱动因素分析与食物网构建}{2025.01 - 2025.02}{}{伦敦}
\role{Excel、文献综述、食物网分析}{}
\begin{itemize}[parsep=0.2ex]
  \item \textbf{数据分析}：基于 ICES 北海 8 种鱼类 90 年长期数据，量化种群动态与物种脆弱性（如鲽形目鱼类因兼捕与晚熟因素在 1990-2000 年间下降约 40\%）
  \item \textbf{文献综合}：结合文献证据解析政策、气候与渔业活动的影响，构建食物网并注释生态过程
  \item \textbf{研究成果}：提出普通海鸦作为指示物种，报告获高分与优秀评价，为渔业资源管理提供科学依据
\end{itemize}

\entry{蜂鸟宏观生态学分析：纬度梯度与生态学规则验证}{2024.10 - 2024.12}{}{伦敦}
\role{R、ggplot2、caper、GLM、PGLS}{}
\begin{itemize}[parsep=0.2ex]
  \item \textbf{假设检验}：基于 AVONET 数据与系统发育树，检验纬度多样性梯度（LDG）、伯格曼法则和拉波波特法则
  \item \textbf{建模分析}：构建物种丰富度热点地图，并应用 GLM 与 PGLS 分析体重、分布范围与纬度关系
  \item \textbf{研究发现}：结果强力支持 LDG；伯格曼法则在普通 GLM 中显著，但在控制系统发育信号（PGLS）后不再显著，揭示迁徙与生态特化的重要作用
\end{itemize}

\entry{环境花色比例对昆虫颜色选择的影响}{2024.09 - 2024.10}{}{开普敦}
\role{野外实验、R（t 检验、Pearson 相关）}{}
\begin{itemize}[parsep=0.2ex]
  \item \textbf{实验设计}：在南非野外课程中，于 6 个样地设计五色陷阱盘实验，研究昆虫访花颜色偏好与环境花色比例的关系
  \item \textbf{实验结果}：黄色、蓝色和白色陷阱均显著高于基线吸引昆虫（p$<$0.001）；黄色捕获率与环境黄花比例显著正相关（r=0.369，p=0.027）
  \item \textbf{研究结论}：支持昆虫利用环境中最常见花色作为觅食线索的假说，深化对授粉生态学的理解
\end{itemize}

\entry{昆虫衰退的驱动因素、影响与保护对策综述}{2024.03 - 2025.06}{}{伦敦}
\role{文献综述、批判性分析}{}
\begin{itemize}[parsep=0.2ex]
  \item \textbf{文献综述}：审阅 70+ 篇文献，系统分析栖息地丧失、农药使用、气候变化等驱动因素，其中一项全球荟萃分析显示昆虫生物量每年下降约 2.5\%
  \item \textbf{影响评估}：评估昆虫衰退对生态系统服务（传粉、养分循环）、农业和人类健康的影响
  \item \textbf{批判分析}：批判性评估“有得有失”的反叙事及地理/分类学报告偏差，提出未来研究与政策方向
\end{itemize}

\entry{温度与基质对海葵胆怯行为的影响}{2024.05 - 2024.06}{}{普利茅斯}
\role{R、Wilcoxon 检验、行为实验}{}
\begin{itemize}[parsep=0.2ex]
  \item \textbf{实验设计}：控制温度与基质条件，测定 52 只海葵的惊愕反应时间（SRT），并进行观察者间信度检验（ICC=0.949）
  \item \textbf{实验结果}：高温显著缩短 SRT（Wilcoxon 检验 p=0.017），表明更“大胆”的行为表现；基质无显著效应
  \item \textbf{研究结论}：研究揭示了环境温度升高与代谢率变化对海洋动物行为的影响
\end{itemize}

\entry{温度暴露时间对蟋蟀攻击行为的影响}{2024.03 - 2024.03}{}{伦敦}
\role{R、GLM（泊松回归）、行为观察}{}
\begin{itemize}[parsep=0.2ex]
  \item \textbf{实验设计}：独立设计实验，测试不同温度暴露时长对地中海田间蟋蟀攻击行为的影响
  \item \textbf{实验结果}：高温显著降低遭遇次数（p=0.040）和持续时间（p=0.013），暴露时间呈负效应趋势
  \item \textbf{研究结论}：研究为理解气候变暖对昆虫行为的即时效应提供证据
\end{itemize}

\entry{生态建模与种群动力学模拟}{2024.01 - 2024.02}{}{伦敦}
\role{R、ggplot2、Lotka-Volterra、逻辑斯蒂增长模型}{}
\begin{itemize}[parsep=0.2ex]
  \item \textbf{Levins集合种群模型}：构建并模拟无限斑块下的种群占领动态，分析 colonization 与 extirpation 的平衡点
  \item \textbf{渔业捕捞模型}：基于逻辑斯蒂增长模型，模拟固定配额与固定努力捕捞策略，计算最大可持续产量（MSY）
  \item \textbf{资源竞争模型}：模拟两种物种对同一资源的竞争过程，验证 R* 理论，分析竞争排斥现象
  \item \textbf{Lotka-Volterra 捕食-被捕食模型}：将两物种模型扩展为三物种捕食链模型，引入随机性与时间滞后，可视化种群周期波动
  \item \textbf{研究成果}：四个建模项目均获得90+评分（最高94分），展现扎实的微分方程建模、算法实现与结果可视化能力，熟练使用R进行动态系统模拟、参数优化与多场景分析
\end{itemize}

\entry{近交系小鼠在医学研究中的应用与局限性综述}{2023.11 - 2023.12}{}{伦敦}
\role{文献综述、遗传学分析}{}
\begin{itemize}[parsep=0.2ex]
  \item \textbf{文献综述}：系统阐述近交小鼠的生成原理、医学应用价值与典型研究案例
  \item \textbf{批判分析}：批判性分析局限性，包括与人类差异、近交衰退及伦理问题
  \item \textbf{研究成果}：完成综述论文并获优秀评价，展示批判性思维与科学写作能力
\end{itemize}

\entry{病毒指纹图谱：利用 HERV-K solo-LTR 多态性进行人群遗传学分析}{2023.11 - 2023.12}{}{伦敦}
\role{PCR、凝胶电泳、群体遗传学分析}{}
\begin{itemize}[parsep=0.2ex]
  \item \textbf{实验操作}：运用 PCR 和凝胶电泳检测 5 组引物对应的基因组多态性，结合 F$_{ST}$ 和 HWE（$\chi^2$ 检验）进行群体遗传学分析
  \item \textbf{分析评估}：评估 HERV-K solo-LTR 作为人群遗传与法医学标记的潜力与挑战
  \item \textbf{研究成果}：完成实验与数据分析，成果获优秀评价
\end{itemize}

\entry{OCA1A 型白化病的生物信息学分析}{2023.10 - 2023.11}{}{伦敦}
\role{NCBI、OMIM、Ensembl、UniProt、AlphaFold、ClinVar}{}
\begin{itemize}[parsep=0.2ex]
  \item \textbf{结构分析}：分析 TYR 基因突变对酪氨酸酶结构与功能的影响，聚焦铜结合域突变，并使用 AlphaFold 结构预测辅助分析
  \item \textbf{数据库整合}：使用多数据库进行变异注释、人群频率分析与跨物种比对
  \item \textbf{研究成果}：完成生物信息学报告，掌握多数据库整合、蛋白结构预测与进化保守性分析技能
\end{itemize}

\section{技能/其他}
\entrySkillsSep
\begin{itemize}[parsep=0.2ex]
  \item \textbf{Python 与开发工具}：Python（主力语言），Git，虚拟环境，Claude Code / Cursor 辅助开发；多后端 LLM 集成（Gemini API、Ollama 本地推理、OpenAI 兼容接口、HuggingFace 多模态模型）、提示工程、JSON Schema 校验、pytest 测试、基于 JSON 配置/状态跟踪的异步子进程编排、基于环境变量的密钥管理
  \item \textbf{Web 开发}：FastAPI（路由/服务架构）、React、TypeScript、Vite、Tailwind CSS
  \item \textbf{分子与生化实验}：PCR 与凝胶电泳（多引物基因分型，如 5 组引物的 solo-LTR 检测）、DNA/RNA 提取（Qiagen 试剂盒）、蛋白质电泳、西方印迹；群体遗传学分析（F$_{ST}$、Hardy-Weinberg 平衡/$\chi^2$ 检验）；跨 NCBI、OMIM、Ensembl、UniProt、ClinVar、AlphaFold 的变异注释与蛋白结构分析
  \item \textbf{生态与野外研究}：野外实验设计（如多样地彩色陷阱盘偏好实验）、具观察者间信度检验的行为学测定（行为学评分、惊愕反应时间测量）、生物多样性/环境调查、宏观生态学与系统发育比较分析（AVONET 级性状数据、PGLS）
  \item \textbf{数据分析与建模}：R（ggplot2、dplyr、caper、ape、vegan、deSolve）、统计推断（GLM/泊松回归、PGLS、Wilcoxon 检验、t 检验、Pearson/Spearman 相关）、基于微分方程与差分方程的种群动力学建模（集合种群、捕食-被捕食、逻辑斯蒂增长与 MSY、资源竞争）
  \item \textbf{语言}：英语（雅思总分7.5），普通话（母语），法语（A2）；\textbf{兴趣爱好}：钢琴，摄影，滑雪
\end{itemize}

\end{document}
